% !TEX root = ../../main.tex
% C12.3 - Future work. 12.3.1 = Vietnamese 9.3 translated (Part I); 12.3.2 = [MERGE] link between the two parts (to be reviewed).
\section{Future Work}
\label{sec:9.3}

\subsection{Part I: The Application}
\label{sec:12.3.1}

The priority is to improve search by text and by attributes. RaSa's re-ranking with its multimodal encoder is now an optional stage of the search flow and is where the model's accuracy lies (Section~\ref{sec:8.4}); making it fast enough for everyday use means computing the image tokens of each appearance at indexing time, in the background worker, instead of at the first search, and running the fusion passes on a GPU, where 128 candidates take well under a second. The measurement of Section~\ref{sec:8.4} also points to a simpler route: a second vector per appearance from a CLIP-like encoder, computed at indexing time and stored in its own Milvus collection, would give text and attribute search close to the re-ranked quality at the cost of one more vector search, while RaSa's vectors keep serving image queries; the design already keeps one collection per encoder version, so the change is contained in the indexing pipeline and the query router. Fine-tuning either model on data from the same domain as surveillance cameras remains the step after that.

For memory, one encoder process shared by the application server and the background worker would remove one of the two copies of the query encoder, about 1~GiB. For speed, inference on the CPU could be optimised, for example with OpenVINO, starting with the image encoder, which takes most of the processing time (Section~\ref{sec:8.5}), or the system could be deployed on a machine with a GPU to analyse many camera streams at the same time instead of in turn. Once the system keeps up with real time, each camera could be processed continuously, and each appearance could be encoded and written as soon as it ends to shorten the delay before it becomes searchable. The write to the three stores already handles each appearance separately, so the change lies mainly in the analysis pipeline.

The evaluation should also compare ByteTrack with BoT-SORT. Further ahead, the system could link the appearances of the same person across cameras, represent each appearance with several frames, and support queries in Vietnamese.

% [MERGE] New paragraph linking the two parts - to be reviewed by both authors.
\subsection{Bringing the Two Parts Together}
\label{sec:12.3.2}

The research of Part~II addresses exactly the weakest point of the application, search by text (Section~\ref{sec:8.4}). Its reasoning layer works on a candidate pool produced by vector retrieval, which is the output the application already obtains from Milvus, so it could be placed behind the application's search step as an optional re-ranking stage. Doing so would require an attribute index built alongside the vectors during indexing, a decision on where the language model runs given the application's on-premises deployment, and a measurement of the added latency per query. The further directions of the reasoning layer itself are stated in Section~\ref{sec:12.2.2}.
